% GENERATED FILE. Edit canonical lessons, then run npm run generate:books.
% canonical-source-hash: fnv1a64:ef8b9bb4ba8760b4
% canonical-lessons: RU-C381-sessiya, RU-C381-zachyot, RU-C381-seminar, RU-C381-konspekt, RU-C381-ball

\chapter{Exam time}
\label{ch:ru-exam-time}
\hlchaptermodality{381}

\begin{chapteropening}
\textbf{By the end of this chapter:} \emph{I can talk about the exam period, a pass-or-fail test, a seminar, my lecture notes and the points I scored.}

Complete the last lesson of chapter 381: I can talk about the exam period, a pass-or-fail test, a seminar, my lecture notes and the points I scored.
\end{chapteropening}

\section[\texorpdfstring{\ru{сессия} (séssiya)}{сессия (séssiya)}]{\ru{сессия} (séssiya) — an exam period (at university); a session}

\label{lesson:RU-C381-sessiya}



\begin{quote}
Before the new one: say the Russian for permanent, then the Russian for temporary.
\end{quote}

\subsection*{You'll want to know: \ru{сессия}}
\textbf{\ru{сессия}} (\emph{séssiya}) — \textquotedblleft{}an exam period (at university); a session\textquotedblright{}.

\begin{quote}
\ru{В} \ru{январе} \ru{у} \ru{студентов} \ru{сессия}.
\end{quote}

\begin{quote}
\emph{V yanvaré u studéntov séssiya.}
\end{quote}

\begin{quote}
In January the students have their exams.
\end{quote}

Feminine. At a Russian university, the weeks of exams at the end of each term. \emph{\ru{Сдавать} \ru{сессию}} — to sit one’s exams.

\begin{cousinweb}
\textbf{\ru{сессия}} is usually traced, through western languages, to Latin \emph{sessiō} \textquotedblleft{}a sitting\textquotedblright{} (Lewis \& Short), built on \emph{sedēre} \textquotedblleft{}to sit\textquotedblright{}: a time when a body sits to do its business.
\end{cousinweb}

\begin{grammarlens}[title={What you've built}]
One more word for telling, explaining and describing at B1.
\end{grammarlens}

\subsection*{Your turn}
\begin{itemize}
\raggedright
  \item \emph{Say it:} \emph{séssiya}
  \item \emph{Say it:} \emph{séssiya}, once more
  \item \emph{Say it:} a sentence of your own with \emph{séssiya}
  \item \emph{Recall:} say the Russian for permanent, then the Russian for temporary, then say \emph{séssiya} again
\end{itemize}

\begin{tcolorbox}[breakable,colback=teal!4,colframe=teal!35!black,title={Before you move on}]
What does \ru{сессия} mean? (\textquotedblleft{}An exam period (at university); a session\textquotedblright{}.) Say it once more.
\end{tcolorbox}
\section[\texorpdfstring{\ru{зачёт} (zachyót)}{зачёт (zachyót)}]{\ru{зачёт} (zachyót) — a pass-or-fail test, a credit}

\label{lesson:RU-C381-zachyot}



\begin{quote}
Before the new one: say the Russian for temporary, then the Russian for an exam period.
\end{quote}

\subsection*{You'll want to know: \ru{зачёт}}
\textbf{\ru{зачёт}} (\emph{zachyót}) — \textquotedblleft{}a pass-or-fail test, a credit\textquotedblright{}.

\begin{quote}
\ru{Завтра} \ru{у} \ru{меня} \ru{зачёт} \ru{по} \ru{истории}.
\end{quote}

\begin{quote}
\emph{Závtra u menyá zachyót po istórii.}
\end{quote}

\begin{quote}
Tomorrow I have a pass-or-fail test in history.
\end{quote}

Masculine. A test graded pass or fail, without a mark. \emph{\ru{Получить} \ru{зачёт}} — to pass it.

\begin{cousinweb}
\textbf{\ru{зачёт}} is built on \emph{\ru{зачесть}} \textquotedblleft{}to count in, to credit\textquotedblright{}, on the old root of \emph{\ru{считать}} and \emph{\ru{число}} \textquotedblleft{}number\textquotedblright{}: a test that is counted in or not.
\end{cousinweb}

\begin{grammarlens}[title={What you've built}]
One more word for telling, explaining and describing at B1.
\end{grammarlens}

\subsection*{Your turn}
\begin{itemize}
\raggedright
  \item \emph{Say it:} \emph{zachyót}
  \item \emph{Say it:} \emph{zachyót}, once more
  \item \emph{Say it:} a sentence of your own with \emph{zachyót}
  \item \emph{Recall:} say the Russian for temporary, then the Russian for an exam period, then say \emph{zachyót} again
\end{itemize}

\begin{tcolorbox}[breakable,colback=teal!4,colframe=teal!35!black,title={Before you move on}]
What does \ru{зачёт} mean? (\textquotedblleft{}A pass-or-fail test, a credit\textquotedblright{}.) Say it once more.
\end{tcolorbox}
\section[\texorpdfstring{\ru{семинар} (seminár)}{семинар (seminár)}]{\ru{семинар} (seminár) — a seminar, a discussion class}

\label{lesson:RU-C381-seminar}



\begin{quote}
Before the new one: say the Russian for an exam period, then the Russian for a pass-or-fail test.
\end{quote}

\subsection*{You'll want to know: \ru{семинар}}
\textbf{\ru{семинар}} (\emph{seminár}) — \textquotedblleft{}a seminar, a discussion class\textquotedblright{}.

\begin{quote}
\ru{На} \ru{семинаре} \ru{мы} \ru{обсуждали} \ru{статью}.
\end{quote}

\begin{quote}
\emph{Na semináre my obsuzhdáli statyú.}
\end{quote}

\begin{quote}
At the seminar we discussed an article.
\end{quote}

Masculine. A class where students talk and answer, as against a \emph{\ru{лекция}}.

\begin{cousinweb}
\textbf{\ru{семинар}} is usually traced, through German \emph{Seminar}, to Latin \emph{sēminārium} \textquotedblleft{}a seed plot, a nursery\textquotedblright{} (Lewis \& Short), built on \emph{sēmen} \textquotedblleft{}seed\textquotedblright{}: a place where young minds are grown.
\end{cousinweb}

\begin{grammarlens}[title={What you've built}]
One more word for telling, explaining and describing at B1.
\end{grammarlens}

\subsection*{Your turn}
\begin{itemize}
\raggedright
  \item \emph{Say it:} \emph{seminár}
  \item \emph{Say it:} \emph{seminár}, once more
  \item \emph{Say it:} a sentence of your own with \emph{seminár}
  \item \emph{Recall:} say the Russian for an exam period, then the Russian for a pass-or-fail test, then say \emph{seminár} again
\end{itemize}

\begin{tcolorbox}[breakable,colback=teal!4,colframe=teal!35!black,title={Before you move on}]
What does \ru{семинар} mean? (\textquotedblleft{}A seminar, a discussion class\textquotedblright{}.) Say it once more.
\end{tcolorbox}
\section[\texorpdfstring{\ru{конспект} (konspékt)}{конспект (konspékt)}]{\ru{конспект} (konspékt) — lecture notes, a summary}

\label{lesson:RU-C381-konspekt}



\begin{quote}
Before the new one: say the Russian for a pass-or-fail test, then the Russian for a seminar.
\end{quote}

\subsection*{You'll want to know: \ru{конспект}}
\textbf{\ru{конспект}} (\emph{konspékt}) — \textquotedblleft{}lecture notes, a summary\textquotedblright{}.

\begin{quote}
\ru{Я} \ru{потерял} \ru{конспект} \ru{лекций}.
\end{quote}

\begin{quote}
\emph{Ya poteryál konspékt léktsiy.}
\end{quote}

\begin{quote}
I have lost my lecture notes.
\end{quote}

Masculine. \emph{\ru{Конспект} \ru{лекций}} — notes from lectures.

\begin{cousinweb}
\textbf{\ru{конспект}} is usually traced to Latin \emph{cōnspectus} \textquotedblleft{}a view, a survey\textquotedblright{} (Lewis \& Short), from \emph{cōnspicere} \textquotedblleft{}to catch sight of\textquotedblright{}, built on \emph{specere} \textquotedblleft{}to see\textquotedblright{}: a summary that lets you take in a subject at a glance.
\end{cousinweb}

\begin{grammarlens}[title={What you've built}]
One more word for telling, explaining and describing at B1.
\end{grammarlens}

\subsection*{Your turn}
\begin{itemize}
\raggedright
  \item \emph{Say it:} \emph{konspékt}
  \item \emph{Say it:} \emph{konspékt}, once more
  \item \emph{Say it:} a sentence of your own with \emph{konspékt}
  \item \emph{Recall:} say the Russian for a pass-or-fail test, then the Russian for a seminar, then say \emph{konspékt} again
\end{itemize}

\begin{tcolorbox}[breakable,colback=teal!4,colframe=teal!35!black,title={Before you move on}]
What does \ru{конспект} mean? (\textquotedblleft{}Lecture notes, a summary\textquotedblright{}.) Say it once more.
\end{tcolorbox}
\section[\texorpdfstring{\ru{балл} (ball)}{балл (ball)}]{\ru{балл} (ball) — a point, a mark (in a test)}

\label{lesson:RU-C381-ball}



\begin{quote}
Before the new one: say the Russian for a seminar, then the Russian for lecture notes.
\end{quote}

\subsection*{You'll want to know: \ru{балл}}
\textbf{\ru{балл}} (\emph{ball}) — \textquotedblleft{}a point, a mark (in a test)\textquotedblright{}.

\begin{quote}
\ru{Он} \ru{получил} \ru{пять} \ru{баллов}.
\end{quote}

\begin{quote}
\emph{On poluchíl pyat bállov.}
\end{quote}

\begin{quote}
He got five points.
\end{quote}

Masculine. \emph{\ru{Средний} \ru{балл}} — a grade average. A dance is \emph{\ru{бал}}, with one \emph{\ru{л}}.

\begin{cousinweb}
\textbf{\ru{балл}} was borrowed from French \emph{balle} \textquotedblleft{}a ball, a bundle\textquotedblright{} (EtymDB); votes and marks were once counted with small balls.
\end{cousinweb}

\begin{grammarlens}[title={What you've built}]
One more word for telling, explaining and describing at B1.
\end{grammarlens}

\subsection*{Your turn}
\begin{itemize}
\raggedright
  \item \emph{Say it:} \emph{ball}
  \item \emph{Say it:} \emph{ball}, once more
  \item \emph{Say it:} a sentence of your own with \emph{ball}
  \item \emph{Recall:} say the Russian for a seminar, then the Russian for lecture notes, then say \emph{ball} again
\end{itemize}

\begin{tcolorbox}[breakable,colback=teal!4,colframe=teal!35!black,title={Before you move on}]
What does \ru{балл} mean? (\textquotedblleft{}A point, a mark (in a test)\textquotedblright{}.) Say it once more.
\end{tcolorbox}
